% Generated draft: requires mathematical and source review.
\documentclass[11pt]{article}
\usepackage[margin=25mm]{geometry}
\usepackage{fontspec}
\setmainfont{Latin Modern Roman}
\usepackage{amsmath,amssymb,mathtools,unicode-math}
\setmathfont{Latin Modern Math}
\usepackage{enumitem,needspace,xurl,hyperref}
\setlength{\parindent}{0pt}
\setlength{\parskip}{5pt}
\setlength{\emergencystretch}{3em}
\newcommand{\N}{\mathbb{N}}
\newcommand{\Z}{\mathbb{Z}}
\newcommand{\Q}{\mathbb{Q}}
\newcommand{\R}{\mathbb{R}}
\newcommand{\C}{\mathbb{C}}
\newcommand{\F}{\mathbb{F}}
\newcommand{\A}{\mathbb{A}}
\newcommand{\PP}{\mathbb{P}}
\newcommand{\E}{\mathbb{E}}
\newcommand{\eps}{\varepsilon}
\newcommand{\dd}{\,\mathrm{d}}
\begin{document}
\section*{UnsolvedMath ID: 30001964}
\textbf{Source problem code:} catalogue entry
\subsection{Definitions and mathematical statement}
The supplied record poses the following question about mathematical practice and its historical reconstruction:
\begin{quote}
How tacit is tacit knowledge in mathematical practice, and how can historians distinguish and recover it from explicit knowledge and incomplete sources?
\end{quote}
The record supplies no operational definitions of tacit knowledge, explicit knowledge, degrees of tacitness, or successful historical recovery. It specifies neither a class of mathematical practices or historical sources nor criteria for distinguishing tacit knowledge from information merely absent from surviving sources. No mathematical objects, hypotheses, boundary conditions, or quantified conclusion are supplied.

The supplied statement-verification note characterizes this as a historical, pedagogical, or methodological discussion prompt rather than an unresolved mathematical proposition. Its registry statement status is \texttt{unrecoverable}, and no separate clean statement is supplied. Consequently, a precise mathematical proposition cannot be recovered from this record without introducing substantive assumptions and an intended target absent from the source data. The question above is retained as the statement-only draft, with these formulation limits explicit.
\subsection{Short English statement}
The question asks how tacit knowledge functions in mathematical practice and how historians can identify and reconstruct it using explicit knowledge and incomplete evidence. The supplied record supports a methodological question but does not specify a mathematical proposition. Its \texttt{open} label is supplied curation, not an independently verified mathematical status.
\subsection{Sources}
\begin{itemize}
\item Local UnsolvedMath registry, supplied record with permanent ID \(30001964\), source problem code \texttt{catalogue entry}, title \emph{Tacit and Explicit Knowledge in Mathematical Practice}; collection: \emph{Oberwolfach Reports Open Problems}. The record gives proposed year 2012 and supplies no named proposer.
\item Supplied source citation: \emph{Explicit Versus Tacit Knowledge in Mathematics} (2012), Oberwolfach Reports, pp.~131--210, DOI 10.4171/owr/2012/04. Supplied source URL and literature-source link, labeled ``original Oberwolfach report'': \url{https://doi.org/10.4171/owr/2012/04}. The supplied background describes Oberwolfach Reports as documenting mathematical workshops and their open-problem discussions.
\item Supplied literature assessment, checked on \texttt{2026-08-21}: status \texttt{open}; the assessment describes a broad conceptual question about tacit knowledge, with no determinate theorem or demonstrably settled answer in the cited literature. This is supplied dated curation derived from the supplied OWR curation, not an independently verified assessment. The listed source checked is the original Oberwolfach report linked above; time-sensitive status and cited claims require independent verification.
\item Supplied statement-verification note: the official history-of-mathematics workshop report treats the item as a historical, pedagogical, or methodological discussion prompt rather than an unresolved mathematical proposition; no separate clean statement is promoted, and the original record remains visible for provenance. This characterization is reproduced from the registry and has not been independently checked.
\end{itemize}
\end{document}
